% Newton's method for one equation: each iterate is the zero of the tangent line.
%
%   figures/tikz/nu-newton-tangent.tex
%       -> media/figures/nu-newton-tangent.{png,svg}
%
% Lecture 15 (Newton Methods for Unconstrained Optimization), Section IV.A.
%
% GEOMETRY (x: 1 unit = 2.3 cm, y: 1 unit = 1.0 cm). f(x) = x^2/2 - 1, f'(x) = x,
% root x* = sqrt(2) = 1.4142. Newton from x^k = 2.6:
%   f(2.6) = 2.38,             tangent zero x^{k+1} = 2.6 - 2.38/2.6 = 1.6846
%   f(1.6846) = 0.4190,        tangent zero x^{k+2} = 1.6846 - 0.4190/1.6846 = 1.4359
% Errors |x - x*| = 1.186, 0.270, 0.0217 (computed in Python, 2026-10-07).
%
% GREYSCALE: f is a thick solid curve, the first tangent solid, the second tangent
% dashed; points are labelled, no distinction is carried by colour alone.
\definecolor{okabeblue}{HTML}{0072B2}

\begin{tikzpicture}[
    x=2.3cm, y=1.0cm,
    >={Latex[length=2.0mm]},
    font=\small,
    lab/.style={font=\scriptsize, inner sep=1.5pt},
  ]
  \draw[black!40, ->] (0.85,0) -- (2.85,0) node[right, lab] {$x$};
  \draw[black!40, ->] (0.95,-0.7) -- (0.95,2.75) node[above, lab] {$f(x)$};
  % f(x) = x^2/2 - 1
  \draw[okabeblue, line width=1.2pt, domain=1.0:2.72, samples=60, variable=\t]
    plot ({\t},{\t*\t/2-1});
  % tangent at x^k = 2.6: y = 2.38 + 2.6 (x - 2.6), zero at 1.6846
  \draw[thick] (1.62,-0.1636) -- (2.72,2.692);
  % tangent at x^{k+1} = 1.6846: y = 0.4190 + 1.6846 (x - 1.6846), zero at 1.4359
  \draw[thick, dashed] (1.38,-0.0941) -- (1.84,0.6807);
  % vertical guides from the axis to the curve
  \draw[densely dotted] (2.6,0) -- (2.6,2.38);
  \draw[densely dotted] (1.6846,0) -- (1.6846,0.4190);
  \filldraw (2.6,2.38) circle (1.5pt);
  \filldraw (1.6846,0.4190) circle (1.5pt);
  \filldraw[fill=white] (2.6,0) circle (1.4pt);
  \filldraw[fill=white] (1.6846,0) circle (1.4pt);
  \filldraw[fill=white] (1.4359,0) circle (1.4pt);
  \filldraw (1.4142,0) circle (1.2pt);
  \node[lab, below] at (2.6,0) {$x^{k}$};
  \node[lab, below] at (1.70,-0.02) {$x^{k+1}$};
  \node[lab, above left] at (1.44,0.02) {$x^{k+2}$};
  \node[lab, below left] at (1.40,-0.05) {$x^{*}$};
  \node[lab, anchor=west] at (2.05,0.75) {tangent at $x^{k}$};
  \node[lab, okabeblue, anchor=east] at (2.15,1.95) {$f(x)$};
\end{tikzpicture}
